\documentclass[11pt]{article}
\usepackage{amsmath,amssymb,amsthm,hyperref}
\newtheorem{theorem}{Theorem}
\newtheorem{lemma}[theorem]{Lemma}
\newtheorem{corollary}[theorem]{Corollary}
\title{An Elementary Rank Bound for Ideal Moment Cancellations}
\author{}
\date{November 23, 2026}
\begin{document}
\maketitle

\paragraph{Scope.}
This is an audit of a proposed number-theoretic construction, not a proof
of the requested universal head bound. It strengthens the special
logarithmic-derivative audit in \texttt{split-product-rank-audit.tex}:
the difference of the products need not split into real affine factors.
All identities in the hypothesis are polynomial identities, not identities
only on the Boolean cube. The proofs below are manual, not Lean checked.

\section{An elementary interpolation fact}

For a finite set $Z\subset\mathbb R^q$, with $q\ge1$, say that $Z$ has the affine
property $\mathrm{CB}(d)$ if every real polynomial of total degree at
most $d$ vanishing on all but one point of $Z$ vanishes on all of $Z$.
This is the affine form of the Cayley--Bacharach condition.

\begin{lemma}[Pairing classes]
Suppose at most $2d$ objects are partitioned into classes, each of size
at most $d$, where $d\ge1$. The objects can be partitioned into at most
$d$ sets, each a singleton or a pair from different classes.
\end{lemma}
\begin{proof}
Add dummy objects in new singleton classes to make the total $2d$.
For exactly $2d$ objects, choose one object from each of the two largest
nonempty classes. There are at least two classes. After removing these
objects, every class has size at most $d-1$: any class of size $d$ was
among the two largest, and there can be no third such class. Induct on
$d$, starting with $d=1$. Finally delete the dummy objects and any
empty pairs. No real pair contains two objects of the same class.
\end{proof}

\begin{theorem}[Small Cayley--Bacharach sets]
Let $d\ge1$ and let $Z$ be a nonempty finite set of distinct real
points having $\mathrm{CB}(d)$.
\begin{enumerate}
\item If $|Z|\le2d+1$, then $Z$ is contained in an affine line.
\item If $|Z|\le2d+2$, then its affine hull has dimension at most two.
\end{enumerate}
\end{theorem}
\begin{proof}
We use the following elementary separation fact: if $p$ is outside the
affine hull of a set $T$, an affine linear polynomial exists which
vanishes on $T$ and is nonzero at $p$. Indeed, after translating $p$
to the origin, the affine hull of $T$ misses zero, so a linear functional
constant and nonzero on that hull gives such a polynomial after
subtracting its constant value. This is finite-dimensional linear
algebra.

Assume first that $Z$ is not collinear and $|Z|\le2d+2$.
No line contains $d+2$ points of $Z$. Otherwise choose $p\in Z$
outside such a line. A hyperplane through the line avoiding $p$,
together with a separate hyperplane through each other point outside
the line and different from $p$, gives a product of at most
$1+(2d+2-(d+2)-1)=d$ affine factors. This product vanishes on
$Z\setminus\{p\}$ and is nonzero at $p$, a contradiction.

Fix $p\in Z$. Partition $Z\setminus\{p\}$ by the lines through $p$.
Each class has at most $d$ points. A pair from different classes has
affine hull avoiding $p$, as does every singleton. If $|Z|\le2d+1$,
the pairing lemma gives at most $d$ such groups. A separating affine
factor for each group again produces a forbidden polynomial. This
proves the first assertion.

For the second assertion it remains to treat $|Z|=2d+2$. Suppose
the affine hull of $Z$ has dimension at least three. The vectors
$z-p$, for $z\in Z\setminus\{p\}$, then span a space of dimension
at least three. We will choose three linearly independent such vectors
whose corresponding points meet every class of size $d$.
For $d=1$, all three vectors are independent and we take them all.
For $d\ge2$, there are at most two classes of size $d$.
If there are two, they account for $2d$ points, leaving just one point.
That last vector must lie outside the span of their two directions.
One vector from each large class and the last vector have the required
property. If there is exactly one large class, extend a vector in its
direction to three independent vectors selected from the spanning set.
If there is none, choose any three independent vectors from that set.

The affine hull of the three chosen points avoids $p$, since their
three difference vectors are linearly independent. After removing
them, there are $2d-2$ points and every remaining line class has size
at most $d-1$. For $d\ge2$ pair them by the lemma; for $d=1$ there
are no remaining points. One factor for the triple and $d-1$ factors
for the pairs give a degree-$d$ polynomial vanishing at every point
except $p$. This is the final contradiction.
\end{proof}

\begin{corollary}[Annihilating moments]
Let $Z\subset\mathbb R^q$ have distinct points and let every weight
$w_z$ be nonzero. Suppose $d\ge1$, $|Z|\le2d+2$, and
\[
 \sum_{z\in Z}w_z P(z)=0
 \quad\text{for every polynomial $P$ of total degree at most $d$}.
\]
Then $Z$ has affine dimension at most two. If $|Z|\le2d+1$, it
is collinear.
\end{corollary}
\begin{proof}
A polynomial vanishing at every point except possibly $p$ satisfies
$w_pP(p)=0$. Since $w_p\ne0$, it also vanishes at $p$. Apply the
theorem.
\end{proof}

\section{Application to affine Prouhet--Tarry--Escott identities}

\begin{theorem}[Three denominator directions]
Let $m\ge1$ and let $u_1,\ldots,u_m,v_1,\ldots,v_m$ be real affine
forms in any number $n$ of variables. Suppose
\[
 F(s,x)=\prod_{j=1}^m(s+u_j(x)),\qquad
 G(s,x)=\prod_{j=1}^m(s+v_j(x)),\qquad F-G=R(x)\ne0
\]
as a polynomial identity. The coefficient vectors of all the $u_j,v_j$
are contained in an affine space of dimension at most two. Consequently,
after substituting any affine form $s=C(x)$, the coefficient vectors
of all $C+u_j,C+v_j$ span a linear space of dimension at most three.
In particular their slope vectors have span dimension at most three.
\end{theorem}
\begin{proof}
The case $m=1$ has at most two coefficient vectors, so is immediate.
Assume $m\ge2$ and put $d=m-1$.
Comparing the coefficients of $s^{m-k}$ gives
$e_k(u_1,\ldots,u_m)=e_k(v_1,\ldots,v_m)$ for $1\le k\le d$.
Newton's identities therefore give the polynomial identities
\[
 \sum_{j=1}^m u_j(x)^k=\sum_{j=1}^m v_j(x)^k
 \quad (0\le k\le d).
\]
For clarity, the induction uses
$p_k-e_1p_{k-1}+e_2p_{k-2}-\cdots+(-1)^k k e_k=0$;
the case $k=0$ simply says that both lists have length $m$.

There can be no identical affine form on opposite sides. Such a form
$u=v$ would make $s+u$ divide $R(x)$ in the polynomial ring, impossible
for a nonzero polynomial independent of $s$.
Let $Z\subset\mathbb R^{n+1}$ be the distinct coefficient vectors,
including the constant coefficient. Give a vector weight equal to its
multiplicity among the $u_j$, or the negative of its multiplicity among
the $v_j$. Thus all weights are nonzero and $|Z|\le2m=2d+2$.

Write a coefficient vector as $z=(z_0,\ldots,z_n)$.
Homogenizing each power-sum identity gives
\[
 \sum_{z\in Z}w_z(z_0t_0+\cdots+z_nt_n)^k=0
 \quad (0\le k\le d).
\]
The coefficient of every degree-$k$ monomial in $t$ is its nonzero
multinomial coefficient times $\sum_z w_z z^\alpha$. Hence all
moments of total degree at most $d$ vanish. Apply the annihilating
moments corollary to obtain an affine plane containing $Z$.

Write its points as $z_\ast+a e+b f$, allowing $e$ or $f$ to be
zero when the affine dimension is smaller. If $c$ is the coefficient
vector of $C$, all substituted denominator vectors lie in
$\operatorname{span}\{c+z_\ast,e,f\}$. Projection onto slope
coordinates cannot increase dimension.
\end{proof}

On a region where all the factors are positive, the score in question is
\[
 \frac{F_s}{F}-\frac{G_s}{G}
 =\sum_j\frac1{s+u_j}-\sum_j\frac1{s+v_j}
 =-\frac{R F_s}{FG},\qquad F_s>0.
\]
It has sign $-\operatorname{sgn}R$ wherever $R\ne0$. The theorem
shows that this construction has denominator slope rank at most three,
without any factorization hypothesis on $R$. Denominator positivity and
the manuscript's common-sign slope restriction must still be checked
separately in an attempted construction.

Apply the bounded-rank-block theorem of
\texttt{split-product-rank-audit.tex} with $r=3$. A sum of
$K=o(H)$ blocks of this form cannot represent every Boolean function
with $H\le C_0 2^n/n^2$ for any fixed $C_0$. This is an obstruction
to a specific ideal-moment construction. It does not exclude unrestricted
head sums, $K$ comparable to $H$, lower-order moment matching, or
identities that hold only on the Boolean cube.

\section{The polynomial-identity hypothesis cannot be dropped}

Here is a concrete boundary example with admissible denominators.
For any $n\ge1$, put $m=n+1$ and define
\[
 u_0(x)=\sum_{i=1}^n x_i,\qquad
 u_k(x)=k(1-x_k)+\sum_{i=k+1}^n x_i\quad(1\le k\le n).
\]
On every Boolean input, these $n+1$ numbers are a permutation of
$0,1,\ldots,n$. This follows by induction: if $x_n=0$, the last
entry is $n$ and the others are the preceding construction; if
$x_n=1$, the last entry is zero and every other entry is increased
by one.

Let $P(s)=\prod_{k=0}^n(s+k)$ and
$\delta=\frac18(\frac34)^n>0$. The polynomial $P(s)-\delta$ has
$n+1$ distinct real roots, one in each interval
$(-k-\frac14,-k+\frac14)$. Indeed the values of $P$ at the two
endpoints have opposite signs and absolute value at least
\[
 \frac14\prod_{j\ne k}\left(|j-k|-\frac14\right)
 \ge\frac14\left(\frac34\right)^n k!(n-k)!
 \ge2\delta.
\]
Subtracting $\delta$ preserves those endpoint signs. The intermediate
value theorem supplies a root in each of the disjoint intervals, and
the degree excludes any further roots. Write the resulting monic
factorization as $P(s)-\delta=\prod_{k=0}^n(s+v_k)$, where each
$v_k$ is a real constant.
Then, on every Boolean input and for every real $s$,
\[
 \prod_{k=0}^n(s+u_k(x))-\prod_{k=0}^n(s+v_k)=\delta.
\]
Nevertheless the slope vectors of the $u_k$ span all $n$ coordinate
directions: the rows for $k=1,\ldots,n$ form an upper triangular
matrix with diagonal entries $-1,-2,\ldots,-n$.

Substitute
$s=C(x)=n+2+(n+2)\sum_i x_i$. Every coefficient of $C+u_k$ and
$C+v_k$ is strictly positive. For the latter constants use
$v_k\in(k-\frac14,k+\frac14)$; for the former slopes use
$n+2-k\ge2$. Since $\sum_{k=0}^n u_k$ is constant as an affine
form, the slope vector of $C$ is the average of those of the
$C+u_k$. The substituted denominators thus still have slope span
dimension $n$.

For $n\ge4$ this contradicts any putative extension of the rank-three
conclusion to Boolean-only identities. It does not contradict the
theorem: off the cube, the displayed product difference need not be
independent of $s$. The example has a constant output sign, so it is
not a compressed construction for arbitrary Boolean labels.

\section{Literature and attribution}

The collinearity statement for at most $2d+1$ points is the affine
counterpart of Bastianelli--Cortini--De Poi, Lemma 2.4,
\emph{The gonality theorem of Noether for hypersurfaces},
\url{https://arxiv.org/pdf/1102.4550}.
Levinson--Ullery, \emph{A Cayley--Bacharach theorem and plane
configurations}, \url{https://arxiv.org/pdf/2102.08525}, study broader
bounds of this kind; see their definition of $\mathrm{CB}(r)$,
Theorems 1.1 and 1.3, and Proposition 2.5.
The elementary proof above establishes the precise real affine statements
used here without importing a general plane-configuration theorem or a
conjectural case of one. The application to ideal moment cancellation is
proved explicitly, rather than inferred from the numerical size of a
Prouhet--Tarry--Escott solution.

\section{Subsequent Boolean-only obstruction}

The full-rank Boolean example above still separates the polynomial
identity rank argument from cube identities. A different invariant
in \texttt{laplace-sign-obstruction.tex} now rules out universality
of Boolean-only ideal product differences at the requested head
rate as well. Their logarithmic derivatives have inverse Laplace
transforms with a common strict sign at all positive times. The
resulting exponential sign-pattern count is too small, irrespective
of denominator slope rank. General rational scores can have
conflicting Laplace signs and are not excluded.

\end{document}
